
\documentclass[11pt,a4paper]{article}

\usepackage[margin=1in]{geometry}
\usepackage[T1]{fontenc}
\usepackage[utf8]{inputenc}
\usepackage{lmodern}
\usepackage{microtype}
\usepackage{amsmath,amssymb}
\usepackage{booktabs}
\usepackage{enumitem}
\usepackage{hyperref}
\usepackage{xcolor}
\usepackage{listings}
\usepackage{fancyhdr}
\usepackage{titlesec}
\usepackage{graphicx}

\hypersetup{
    colorlinks=true,
    linkcolor=blue,
    urlcolor=blue,
    citecolor=blue
}

\definecolor{codegray}{gray}{0.95}

\lstset{
    basicstyle=\ttfamily\small,
    backgroundcolor=\color{codegray},
    frame=single,
    breaklines=true,
    columns=fullflexible,
    keepspaces=true,
    showstringspaces=false
}

\pagestyle{fancy}
\fancyhf{}
\lhead{FRC Research Notes}
\rhead{Week 4}
\cfoot{\thepage}

\titleformat{\section}{\large\bfseries}{\thesection.}{0.5em}{}
\titleformat{\subsection}{\normalsize\bfseries}{\thesubsection.}{0.5em}{}

\title{\textbf{FRC Research --- Week 4 Working Notes}\\
\large 4~$\mu$s Formation Extension and Transition to GS/MRR Tilt Reproduction}
\author{Juan Pablo Solís Ruiz\\
Supervisor: Prof. Ren Haijun}
\date{September 28--October 4, 2026}

\begin{document}

\maketitle
\tableofcontents
\newpage

% ============================================================
\section{Purpose of These Notes}

These are the working research notes for Week 4 of the FRC project.

Week 3 completed the software-takeover and inherited-baseline reproduction
stage. The original 2000-step WarpX formation case was parameterized,
generated through Framework v0.1, executed on SCNet, and quantitatively
compared against the two frozen legacy realizations.

The principal result carried into this week is

\[
\boxed{
\text{Framework v0.1 inherited-baseline regression: PASS}
}
\]

for the macroscopic magnetic-field and density evolution of the inherited
formation case.

The purpose of Week 4 is therefore to:

\begin{enumerate}
    \item perform one controlled extension of the validated formation case
    from \(2~\mu\mathrm{s}\) to the full imposed-field rise time
    \(t_{\mathrm{peak}}=4~\mu\mathrm{s}\);
    \item characterize the additional \(2\rightarrow4~\mu\mathrm{s}\)
    evolution without reopening the completed baseline-regression work;
    \item freeze formation as a validated framework example;
    \item move the main research effort to the Grad--Shafranov / Modified
    Rigid Rotor equilibrium and three-dimensional tilt-instability problem.
\end{enumerate}

% ============================================================
\section{Handoff from Week 3}

At the end of Week 3, the inherited 2000-step formation case had been
reproduced using Framework v0.1.

The validated reference configuration has

\[
N_t = 2000,
\qquad
\Delta t = 1~\mathrm{ns},
\qquad
t_{\mathrm{end}}=2~\mu\mathrm{s},
\]

while the imposed external-field ramp uses

\[
t_{\mathrm{peak}} = 4~\mu\mathrm{s}.
\]

Thus, the inherited case terminates when the drive amplitude is only

\[
A_{\mathrm{time}}(2~\mu\mathrm{s})
=
\sin\left(
\frac{\pi(2~\mu\mathrm{s})}
     {2(4~\mu\mathrm{s})}
\right)
=
\frac{1}{\sqrt{2}}
\simeq 0.7071.
\]

This motivated one final formation run extending the same validated model to

\[
t=4~\mu\mathrm{s},
\]

where

\[
A_{\mathrm{time}}=1.
\]

The purpose of this extension is exploratory rather than regression-based:
there is no frozen 4000-step legacy run against which to compare it.

% ============================================================
\section{Research Log --- 2026-09-28}

\subsection{Decision to Extend the Formation Case to 4~\texorpdfstring{$\mu$s}{us}}

The 2000-step baseline was intentionally left unchanged and retained as the
validated regression reference.

A separate case configuration was derived for the longer simulation:

\begin{lstlisting}
configs/formation/formation_4us.toml
\end{lstlisting}

with the intended case identity

\begin{lstlisting}
name = "formation_4us_ramp_complete_4000"
purpose = "formation_extension"
max_steps = 4000
\end{lstlisting}

The timestep was not changed:

\[
\Delta t = 10^{-9}\ \mathrm{s},
\]

so

\[
t_{\mathrm{end}}
=
4000\Delta t
=
4.0\times10^{-6}\ \mathrm{s}
=
4~\mu\mathrm{s}.
\]

No physical model, geometry, particle loading, transport model, field-drive
formula, or boundary condition was intentionally changed relative to the
validated baseline.

The longer run therefore isolates the effect of continuing the already
validated formation simulation through the remainder of the imposed-field
rise.

\subsection{Output Strategy}

The plotfile diagnostic cadence was retained so that the full evolution could
be inspected every 100 steps.

Checkpoint interval

\[
N_{\mathrm{chk}} = 2000
\]

was also retained, producing checkpoints at approximately

\[
t=0,\ 2,\ 4~\mu\mathrm{s}.
\]

The checkpoint at step 2000 provides a direct connection to the endpoint of
the previously validated baseline, while the final checkpoint preserves the
new \(4~\mu\mathrm{s}\) state for possible later restart or analysis.

\subsection{SCNet Execution}

The formation extension was submitted through Framework v0.1 as Slurm job

\begin{lstlisting}
123207160
\end{lstlisting}

and executed on

\begin{lstlisting}
a02r2n15
\end{lstlisting}

with the same CPU layout used by the validated baseline:

\[
4~\mathrm{MPI~ranks}
\times
8~\mathrm{OpenMP~threads/rank}
=
32~\mathrm{allocated~CPU~cores}.
\]

The completed accounting record was

\begin{lstlisting}
State    = COMPLETED
Elapsed  = 05:29:56
ExitCode = 0:0
\end{lstlisting}

The run produced approximately

\[
18~\mathrm{GB}
\]

of output, including

\[
41
\]

plotfile diagnostics and

\[
3
\]

checkpoint directories.

This is consistent with diagnostics from step 0 through step 4000 in
increments of 100 and checkpoints at steps 0, 2000, and 4000.

% ============================================================
\section{Research Log --- 2026-09-29}

\subsection{Quick Analysis Objective}

The purpose of the analysis was not to repeat the full Week 3 regression.

Instead, the analysis focused on the new interval

\[
2~\mu\mathrm{s}
\leq t
\leq
4~\mu\mathrm{s},
\]

while retaining the earlier evolution as context.

The main quantities inspected were

\[
B_{z,\mathrm{center}},
\qquad
B_{z,\min},
\qquad
B_{z,\max},
\]

\[
\rho_{\mathrm{center}},
\qquad
\rho_{\max},
\]

and

\[
T_{i,\mathrm{center}},
\qquad
T_{i,\max}.
\]

Spatial \(x-z\) midplane slices and radial \(B_z\) profiles were also generated
at

\[
t=2,\ 3,\ 4~\mu\mathrm{s}.
\]

As in the earlier baseline audit, \(B_z=0\) is used only as a simple
field-reversal indicator. It is \textbf{not} treated as a formally reconstructed
magnetic separatrix.

\subsection{Selected Scalar States}

Table~\ref{tab:formation4us_selected} summarizes representative times from the
full \(0\rightarrow4~\mu\mathrm{s}\) evolution.

\begin{table}[htbp]
    \centering
    \small
    \caption{Selected states of the 4~$\mu$s formation extension.}
    \label{tab:formation4us_selected}
    \begin{tabular}{@{}rrrrrr@{}}
        \toprule
        \(t\) [$\mu$s] &
        \(A_{\mathrm{time}}\) &
        \(B_{z,c}\) [T] &
        \(B_{z,\min}\) [T] &
        \(\rho_c\) &
        \(\rho_{\max}\) \\
        \midrule
        0.0 & 0.000 & -0.0800 & -0.0800 &  8.01 &  8.07 \\
        1.0 & 0.383 & -0.0801 & -0.1414 &  7.93 & 21.64 \\
        1.7 & 0.619 & -0.3947 & -0.4831 & 36.24 & 44.15 \\
        2.0 & 0.707 & -0.1240 & -0.2507 & 12.67 & 50.82 \\
        3.0 & 0.924 & -0.1464 & -0.1982 & 12.86 & 68.29 \\
        4.0 & 1.000 & -0.3093 & -0.3103 & 24.03 & 80.19 \\
        \bottomrule
    \end{tabular}
\end{table}

The most negative central magnetic field over the complete run occurred at

\[
t\simeq1.7~\mu\mathrm{s},
\]

with

\[
B_{z,\mathrm{center}}
\simeq
-0.395~\mathrm{T}.
\]

The most negative chamber field also occurred at approximately the same time:

\[
B_{z,\min}
\simeq
-0.483~\mathrm{T}.
\]

This confirms that the strongest early magnetic compression does
\textbf{not} occur at the end of the imposed-field ramp.

\subsection{Non-Monotonic Magnetic Response}

The imposed drive rises monotonically from

\[
A_{\mathrm{time}}=0
\]

to

\[
A_{\mathrm{time}}=1
\]

at \(4~\mu\mathrm{s}\).

The plasma response is not monotonic.

The central field evolves approximately as

\[
-0.080~\mathrm{T}
\rightarrow
-0.395~\mathrm{T}
\rightarrow
-0.124~\mathrm{T}
\rightarrow
-0.309~\mathrm{T},
\]

at

\[
t=
0,\ 1.7,\ 2.0,\ 4.0~\mu\mathrm{s},
\]

respectively.

The current working interpretation is therefore

\[
\boxed{
\text{strong compression}
\rightarrow
\text{relaxation}
\rightarrow
\text{renewed compression}
}
\]

rather than a quasi-static monotonic response that simply follows the imposed
field ramp.

The words ``relaxation'' and ``renewed compression'' are presently
descriptive. A more rigorous dynamical interpretation would require additional
energy, flux, force-balance, and topology diagnostics.

\subsection{Magnetic Structure from 2 to 4~\texorpdfstring{$\mu$s}{us}}

Figure~\ref{fig:bz_2_3_4us} shows the midplane axial-field structure at the
validated baseline endpoint and at the two additional times analyzed during
the extension.

\begin{figure}[htbp]
    \centering

    \begin{minipage}{0.32\textwidth}
        \centering
        \includegraphics[
            width=\linewidth
        ]{img/formation_4us/Bz_xz_step002000.png}

        \small \(t=2~\mu\mathrm{s}\)
    \end{minipage}
    \hfill
    \begin{minipage}{0.32\textwidth}
        \centering
        \includegraphics[
            width=\linewidth
        ]{img/formation_4us/Bz_xz_step003000.png}

        \small \(t=3~\mu\mathrm{s}\)
    \end{minipage}
    \hfill
    \begin{minipage}{0.32\textwidth}
        \centering
        \includegraphics[
            width=\linewidth
        ]{img/formation_4us/Bz_xz_step004000.png}

        \small \(t=4~\mu\mathrm{s}\)
    \end{minipage}

    \caption{
        Midplane \(B_z(x,z)\) evolution during the second half of the
        imposed-field rise. The black \(B_z=0\) contour, when present, is used
        only as a field-reversal proxy and not as a formal magnetic separatrix.
    }
    \label{fig:bz_2_3_4us}
\end{figure}

The simulation retains regions of both negative and positive axial field at
the end of the field ramp.

At \(t=4~\mu\mathrm{s}\),

\[
B_{z,\mathrm{center}}
\simeq
-0.309~\mathrm{T},
\]

while

\[
B_{z,\max}
\simeq
+0.245~\mathrm{T}.
\]

Therefore, field reversal remains present when

\[
A_{\mathrm{time}}=1.
\]

This result demonstrates persistence of a reversed-field structure through
the complete imposed-field rise, but it is not by itself sufficient to
establish the exact magnetic separatrix, O-point, X-point, or closed-flux
volume.

Those quantities will eventually require reconstruction of a poloidal flux
function such as \(\psi(r,z)\).

\subsection{Radial Magnetic Evolution}

Figure~\ref{fig:bz_radial_4us} compares the radial midplane \(B_z\) profiles
at 2, 3, and \(4~\mu\mathrm{s}\).

\begin{figure}[htbp]
    \centering
    \includegraphics[
        width=0.82\textwidth
    ]{img/formation_4us/Bz_radial_2_3_4us.png}

    \caption{
        Radial \(B_z\) profile through the midplane at
        \(t=2,\ 3,\ 4~\mu\mathrm{s}\).
        The evolving zero crossings and profile amplitudes show that the
        magnetic response continues to reorganize during the second half of
        the imposed-field ramp.
    }
    \label{fig:bz_radial_4us}
\end{figure}

The radial evolution supports the scalar observation that the magnetic
configuration does not merely scale uniformly with the imposed field.

Instead, the field structure changes dynamically as the ramp proceeds.

\subsection{Density Evolution}

The chamber maximum density continues to increase after the original
\(2~\mu\mathrm{s}\) baseline endpoint:

\[
\rho_{\max}(2~\mu\mathrm{s})
\simeq 50.8,
\]

\[
\rho_{\max}(3~\mu\mathrm{s})
\simeq 68.3,
\]

and

\[
\rho_{\max}(4~\mu\mathrm{s})
\simeq 80.2.
\]

The largest sampled chamber maximum occurs slightly before the end of the run,

\[
\rho_{\max}
\simeq
83.65
\]

at approximately

\[
t=3.8~\mu\mathrm{s}.
\]

The central density does not increase in the same monotonic way.

It evolves from approximately

\[
\rho_{\mathrm{center}}
\simeq 36.2
\]

at the strong early compression near \(1.7~\mu\mathrm{s}\) to

\[
12.7
\]

at \(2~\mu\mathrm{s}\), then rises again to approximately

\[
24.0
\]

at \(4~\mu\mathrm{s}\).

Therefore, the late increase of \(\rho_{\max}\) should not be interpreted as
uniform central compression.

A future detailed study should locate the region responsible for the density
maximum and relate it to the flux topology and plasma boundary.

\subsection{Ion-Temperature Diagnostics}

The central ion-temperature diagnostic follows the strong early compression
more closely than the chamber maximum.

The largest central value occurs around

\[
t=1.7~\mu\mathrm{s},
\]

with

\[
T_{i,\mathrm{center}}
\simeq 23.2
\]

in the units reported by the current WarpX/yt diagnostic pipeline.

By contrast, the pointwise chamber maximum reaches values above

\[
3\times10^3,
\]

with the sampled maximum

\[
T_{i,\max}
\simeq
3.41\times10^3
\]

near

\[
t=3.7~\mu\mathrm{s}.
\]

This quantity remains unsuitable as a strict macroscopic physical diagnostic.

The earlier regression already demonstrated that pointwise
\(T_{i,\max}\) is highly sensitive to stochastic PIC sampling.

In addition, the exact physical-unit interpretation of the temperature field
should be independently verified before interpreting the large extrema as a
quantitative plasma temperature.

\subsection{Important Extrema over 0--4~\texorpdfstring{$\mu$s}{us}}

The automatic diagnostic scan found:

\begin{table}[htbp]
    \centering
    \caption{Selected extrema in the 4000-step formation extension.}
    \begin{tabular}{@{}lll@{}}
        \toprule
        Quantity & Extremum & Time \\
        \midrule
        \(B_{z,\mathrm{center}}\) &
        \(-3.947\times10^{-1}\) T &
        \(1.7~\mu\mathrm{s}\) \\
        \(B_{z,\min}\) &
        \(-4.831\times10^{-1}\) T &
        \(1.7~\mu\mathrm{s}\) \\
        \(B_{z,\max}\) &
        \(+2.600\times10^{-1}\) T &
        \(2.6~\mu\mathrm{s}\) \\
        \(\rho_{\max}\) &
        \(8.365\times10^{1}\) &
        \(3.8~\mu\mathrm{s}\) \\
        \(T_{i,\mathrm{center}}\) &
        \(2.322\times10^{1}\) &
        \(1.7~\mu\mathrm{s}\) \\
        \(T_{i,\max}\) &
        \(3.415\times10^{3}\) &
        \(3.7~\mu\mathrm{s}\) \\
        \bottomrule
    \end{tabular}
\end{table}

The fact that these extrema occur at different times is another indication
that the system has significant transient dynamics and cannot be characterized
by a single scalar ``compression factor'' tied directly to the external ramp.

% ============================================================
\section{Working Interpretation of the 4~\texorpdfstring{$\mu$s}{us} Extension}

The main observations from the extension are:

\begin{itemize}
    \item The strongest sampled central magnetic compression occurs near
    \(1.7~\mu\mathrm{s}\), well before completion of the imposed-field ramp.
    \item A strong relaxation follows before \(2~\mu\mathrm{s}\).
    \item The magnetic field subsequently develops a second compression /
    reorganization stage as the drive approaches its final value.
    \item Field reversal remains present at \(4~\mu\mathrm{s}\).
    \item The chamber maximum density continues increasing until approximately
    \(3.8~\mu\mathrm{s}\), while the central density behaves non-monotonically.
    \item Large pointwise \(T_i\) maxima are treated as noisy diagnostics rather
    than representative bulk plasma temperatures.
\end{itemize}

These observations are sufficient for the present purpose of the formation
extension.

A rigorous formation/compression study would eventually require additional
quantities, including

\[
\psi(r,z),
\quad
R_s,
\quad
L_s,
\quad
\Phi_{\mathrm{trapped}},
\quad
E_{\mathrm{mag}},
\quad
E_{\mathrm{kin}},
\]

as well as particle and energy-loss diagnostics.

Those additions remain relevant to the long-term framework, but they are not
required before moving to the next short-term research objective.

% ============================================================
\section{Formation Stage Status}

The formation workflow has now passed through the following sequence:

\[
\text{legacy takeover}
\rightarrow
\text{baseline audit}
\rightarrow
\text{Framework v0.1}
\rightarrow
\text{smoke test}
\]

\[
\rightarrow
\text{2000-step regression}
\rightarrow
\text{4~$\mu$s extension}.
\]

The validated 2000-step case remains the primary regression reference.

The 4000-step case is retained as the first controlled scientific extension of
that reference.

At this point, additional long formation runs are not the highest-priority use
of compute time.

The formation module should instead be revisited when new diagnostics,
parameter scans, or experimental questions require it.

% ============================================================
\section{Next Research Stage: GS/MRR and Tilt Reproduction}

The main short-term project objective now moves to the inherited
Grad--Shafranov / Modified Rigid Rotor equilibrium workflow and the
three-dimensional tilt-instability simulations.

The intended sequence is

\[
\boxed{
\text{GS/MRR equilibrium}
\rightarrow
\text{WarpX initialization}
\rightarrow
\text{legacy tilt reproduction}
\rightarrow
\text{validated tilt diagnostics}
}
\]

before introducing new parameter studies.

The known audit questions to resolve include:

\begin{enumerate}
    \item \textbf{Alfv\'en-time normalization}

    The inherited material contains apparently inconsistent values for the
    Alfv\'en time. The normalization used in reported growth rates must be
    reconstructed from the actual equilibrium parameters.

    \item \textbf{Ion skin depth and \(S^*\)}

    The current implementation must be checked for consistency between
    deuterium used in the simulation and the mass used when calculating

    \[
    d_i = \frac{c}{\omega_{pi}}.
    \]

    A proton/deuterium mismatch would modify \(d_i\) and therefore \(S^*\).

    \item \textbf{Tilt-mode diagnostic}

    The inherited approximation relating \(v_z\), \(j_z\), and
    \(\rho_{\mathrm{protons}}\) must be audited.

    A direct particle bulk-velocity diagnostic may be required to avoid
    confusing total current with ion flow.

    \item \textbf{Growth-rate extraction}

    The procedure used to obtain

    \[
    \gamma_{\mathrm{tilt}}
    \]

    and the normalization

    \[
    \gamma_{\mathrm{tilt}}\tau_A
    \]

    must be reconstructed and tested quantitatively.

    \item \textbf{Seeded and unseeded runs}

    The inherited \(n=1\) tilt seed should be reproduced explicitly and
    compared against the unseeded case before modifying the equilibrium or
    stability parameters.
\end{enumerate}

% ============================================================
\section{Framework Development Direction for Week 4}

Framework v0.1 has now demonstrated that it can safely reproduce and extend the
formation workflow.

The next framework growth should be driven by the GS/MRR and tilt problem
rather than by adding generic infrastructure in advance.

Likely next components are:

\begin{itemize}
    \item a structured equilibrium configuration;
    \item provenance for the generated \texttt{init.bp} openPMD file;
    \item explicit equilibrium-derived quantities such as
    \(R_0\), \(Z_0\), \(V_A\), \(\tau_A\), \(d_i\), and \(S^*\);
    \item tilt-seed configuration as a first-class parameter;
    \item reproducible extraction of the \(n=1\) mode amplitude;
    \item growth-rate fitting and normalization;
    \item comparison tools for seeded/unseeded and field-strength cases.
\end{itemize}

The same project principle remains in force:

\[
\boxed{
\text{reproduce first, then change physics}
}
\]

No refactor or model change should be mixed with the first GS/MRR tilt
reproduction.

% ============================================================
\section{Week 4 Working Status}

As of the end of 2026-09-29:

\begin{itemize}
    \item the inherited 2000-step formation baseline remains validated;
    \item Framework v0.1 has completed both smoke and full regression tests;
    \item the same validated formation model has been extended successfully to
    \(4~\mu\mathrm{s}\);
    \item the 4000-step job completed in \(05{:}29{:}56\) with exit code
    \(0{:}0\);
    \item the extended run produced 41 diagnostics, 3 checkpoints, and
    approximately 18~GB of output;
    \item the strongest early compression occurs around
    \(1.7~\mu\mathrm{s}\);
    \item a subsequent relaxation and renewed magnetic compression are observed;
    \item field reversal persists through the end of the imposed-field ramp;
    \item formation is now sufficiently characterized for the project to move
    to GS/MRR and tilt-instability reproduction.
\end{itemize}

\[
\boxed{
\text{Next: reproduce and audit the inherited GS/MRR tilt workflow}
}
\]

\end{document}
